\section{Related Work}
\label{sec:related}

This section reviews the evolution of fault diagnosis methodologies and identifies the critical gaps our proposed model aims to address. We first briefly discuss traditional fault diagnosis techniques, highlighting their limitations in modern data-driven environments. Subsequently, we explore the advancements and current challenges in zero-shot fault diagnosis and generative diffusion models. Finally, we examine the recent emergence of Physics-Informed Neural Networks (PINN) and their potential to enforce physical viability in generative tasks.

\subsection{Traditional Fault Diagnosis}
Traditional fault diagnosis in rotating machinery has historically relied on analytical modeling, signal processing techniques, and expert systems \cite{randall2011rolling,jardine2006review}. Analytical models utilize profound dynamic equations to monitor structural health, but they are often too complex to accurately formulate for intricate, nonlinear industrial systems \cite{jardine2006review}. Conversely, signal processing techniques, such as wavelet transforms and empirical mode decomposition, require extensive domain expertise to extract discriminative features from raw vibration data \cite{yan2014wavelets}. While expert systems successfully integrate these qualitative experiences and quantitative features, they generally lack the scalability and diagnostic flexibility required for highly variable operating conditions \cite{zhao2019deep}. Consequently, deep learning approaches have largely superseded analytical and expert-driven traditional methods, though they remain heavily dependent on massive, comprehensively annotated fault datasets \cite{lei2020applications}.

\subsection{Zero-Shot Fault Diagnosis}
Zero-Shot Learning (ZSL) effectively mitigates the data scarcity challenge by establishing a cognitive bridge between seen classes and unseen classes via semantic mapping \cite{xian2018zero}. In rotating machinery diagnostics, early attempts largely relied on attribute-based semantic embeddings mapping multi-domain features to fault descriptions \cite{zhang2023zero}. More recently, semantic approaches have broadened to include NLP embeddings and diffusion-encoded probability \cite{chen2024semantic}. Recent work has explored knowledge distillation-based ZSL frameworks \cite{liu2025knowledge} and generalized ZSL for compound faults \cite{song2025generalized}, demonstrating improved performance in recognizing unseen fault categories. However, conventional mapping-based ZSL models often suffer from the hubness problem and poor generalization \cite{xian2018zero}. Generative zero-shot approaches emerged to supplement this by synthesizing hypothetical data for the unseen fault classes, converting ZSL into fully supervised contexts \cite{schonfeld2019cada,kingma2013auto}. Despite advancements, most existing models perform deterministic transformations completely disregarding the intrinsic uncertainty inherent in signal measurement and generative extrapolations \cite{gal2016dropout}.

\subsection{Generative Diffusion Models}
Diffusion models operate by progressively destroying signal structures with Gaussian noise and subsequently learning to reverse this corruption process \cite{ho2020denoising}. In domains such as image and audio generation, diffusion models have eclipsed VAEs and GANs concerning generation quality and training stability \cite{song2020score,goodfellow2014generative}. Recent literature exhibits preliminary implementations of diffusion architectures for bearing fault augmentation \cite{ho2020denoising}. Nonetheless, when directly deployed on time-series vibration domains, naive diffusion models blindly replicate statistical noise and generate unphysical artifacts \cite{song2020score}. This leaves a significant gap wherein generated temporal data lacks adherence to dynamic equations required for industrial simulation-grade virtual samples.

\subsection{Large Language Models in Fault Diagnosis}
The emergence of Large Language Models (LLMs) has opened new avenues for fault diagnosis by leveraging their semantic understanding and reasoning capabilities. Recent work includes FD-LLM \cite{fdllm2025}, which adapts LLMs to handle numerical sensor data for identifying machine faults, and multimodal LLM frameworks \cite{mllm2024} that combine visual and textual information for fault detection. LLMs have demonstrated strong adaptability across different operational conditions \cite{arxiv2024fdllm}, outperforming traditional deep learning approaches in certain scenarios. However, the integration of LLM-generated semantic embeddings with physics-constrained generative models remains largely unexplored.

\subsection{Physics-Informed Neural Networks (PINN)}
Physics-Informed Neural Networks systematically integrate known differential equations into the deep learning optimization cycle \cite{raissi2019physics}. Serving as a hybrid of data-driven and mechanistic principles, PINNs have shown exceptional capability in surrogate modeling of fluid dynamics and thermodynamic state estimation \cite{chen2021physics}. Recent applications in fault diagnosis include multistage PINN for valve fault detection \cite{energies2024multistage}, inverse PINN for imbalanced datasets \cite{plosone2025inverse}, and online PINN for remaining useful life prediction \cite{nature2025online}. In intelligent fault diagnosis, incorporating physics restricts the hypothesis space, thereby mitigating overfitting and ensuring theoretical feasibility \cite{yucesan2022hybrid}. Exploiting physical constraints as a strict regularization boundary for generative models in the zero-shot learning domain remains thoroughly unexplored.
